\section{Production Deployment}
\label{sec:production}

\subsection{Q-Chain finality on Lux primary network}

The current production deployment of QuasarCert is on the Lux
primary network's Q-Chain, where it provides PQ-safe finality for
the cross-chain Warp messaging layer. The Q-Chain runs the
\texttt{Pulsar} cert profile (BLS + Pulsar + Z-Chain Groth16), which
is the smallest profile that satisfies the PQ-floor requirement.

Q-Chain validators are the Lux primary network validators (the same
validator set that secures the P-Chain, C-Chain, and X-Chain). Each
validator holds:

\begin{itemize}[leftmargin=2em,nosep]
\item A BLS12-381 keypair with PoP for the QuasarCert BLS leg.
\item A Pulsar threshold-secret share for the QuasarCert Pulsar leg.
\item A per-validator ML-DSA-65 keypair for the Z-Chain Groth16 leg.
\end{itemize}

The Pulsar DKG is run once per validator-set epoch by Q-Chain
validators; the resulting group public key $\mathit{pk}^{\mathrm{Puls}}$
is committed to the validator-set Merkle root for that epoch.

\subsection{LP-020 Quasar consensus activation}

The Quasar 3.0 protocol was activated on the Lux primary network on
$2025$-$12$-$25$ per LP-020 \cite{lp020}. The activation enabled the
three-lane cert (BLS + Pulsar + Z-Chain Groth16) on Q-Chain blocks
from height $h_{\mathrm{act}}$ onwards. Pre-activation blocks carry
BLS-only certs (legacy); post-activation blocks carry full
QuasarCerts under the Pulsar profile.

The activation roadmap (LP-020 \S14) committed to:

\begin{itemize}[leftmargin=2em,nosep]
\item Pulsar primitive at Tier A submission readiness (achieved
      $2026$-$05$-$18$ per the \texttt{CRYPTOGRAPHER-SIGN-OFF.md}
      in \texttt{luxfi/pulsar} commit \texttt{443c725}).
\item Aurora profile reserved for the deployment phase after
      Corona's Tier A audit (target: late 2026, gated on Corona's
      open GATE-1/2/3/4 closure per \texttt{luxfi/corona}'s
      submission status).
\item Polaris profile reserved for after Magnetar's Tier A audit
      (target: 2027 Q3 per \texttt{luxfi/magnetar/SUBMISSION-STATUS.md}).
\end{itemize}

\subsection{Profile gating for Aurora and Polaris}

The Aurora and Polaris profiles are operationally ready in the
\texttt{luxfi/consensus} implementation but are not yet enabled by
default on any Lux production chain. The gating is documented in
\texttt{luxfi/consensus/config/pq\_mode.go} and is bound to the
maturity tier of the non-Pulsar primitives.

\paragraph{Aurora gating.}
Aurora requires Corona at Tier A submission readiness. As of
$2026$-$05$-$18$, Corona is at ``Tier A documentation shape
complete'' with four open gates: EasyCrypt theory shells (GATE-1,
v0.7.0), Lean$\leftrightarrow$EasyCrypt bridge (GATE-2, v0.7.0),
dudect $10^9$ on threshold layer (GATE-3, v0.8.0), and external
audit (GATE-4, v0.8.0). The Aurora profile is unblocked once GATE-4
closes.

\paragraph{Polaris gating.}
Polaris requires both Corona Tier A and Magnetar Tier A. Magnetar
is currently at Tier B (v0.1 reveal-and-aggregate implementation
landed with internal review; full Tier A artifact bundle is on the
roadmap). Polaris activation target is post-Magnetar-Tier-A
(2027 Q3).

The gating discipline is: a profile is enabled in production only
after every primitive on the profile has external audit and
mechanized proofs (EasyCrypt / Lean / Jasmin). This is the
production-grade bar; primitives at lower tiers are research
references, not deployment surfaces.

\subsection{Cross-chain finality propagation}

The QuasarCert is the finality envelope for the Lux ecosystem's
cross-chain messaging \cite{lp075}. A message from chain $c_1$ to
chain $c_2$ carries the source chain's QuasarCert; the destination
chain verifies the cert against the source chain's known public
keys.

The wire-format stability commitment (Section~\ref{sec:wire})
is what makes this work: a Q-Chain QuasarCert, a C-Chain
QuasarCert, and a chain-VM EVM QuasarCert all use the same Go struct,
the same encoder, and the same verifier. The only variation across
chains is the cert profile (\texttt{Pulsar} for Q-Chain;
potentially \texttt{Aurora} or \texttt{Polaris} for high-assurance
L1 chains in the future) and the validator-set / primitive public
keys.

A Lux chain that adopts the Aurora profile can produce QuasarCerts
that are interoperable with Pulsar-profile chains via downward
compatibility: an Aurora cert verifies on a Pulsar-only verifier
that ignores the Corona leg, with the verification rule degraded to
profile-Pulsar correctness. (In practice, Lux disallows this
downward verification because the profile-match clause of
Definition~\ref{def:verify} would reject the cert as malformed for
the wrong profile. Cross-profile interoperability requires both
chains to be configured for the same profile or for a profile
rotation event.)

\subsection{Operational runbooks}

Each primitive ships an operational runbook documenting its
deployment, key-rotation, and incident-response procedures.

\begin{itemize}[leftmargin=2em,nosep]
\item Pulsar: \texttt{luxfi/pulsar/DEPLOYMENT-RUNBOOK.md} covers DKG
      coordination, group-key publication, resharing across
      validator-set rotations, identifiable abort handling, and
      the v0.1 aggregator-TCB caveat with mitigation guidance.
\item Corona: \texttt{luxfi/corona/DEPLOYMENT-RUNBOOK.md} (Tier A
      docs landed; equivalent shape to Pulsar's runbook).
\item Magnetar: \texttt{luxfi/magnetar/DEPLOYMENT-RUNBOOK.md} covers
      the byte-wise Shamir VSS DKG, the reveal-and-aggregate signing
      flow, and the aggregator zeroization protocol.
\item Quasar umbrella: profile rotation, primitive version bumps,
      and cross-repo cascade procedures documented at
      \texttt{luxfi/quasar/SPEC.md} \S Wire-format stability commitment.
\end{itemize}

\subsection{NIST MPTC submission trajectory}

Each PQ primitive in QuasarCert is being submitted to the NIST
MPTC project \cite{nistir8214c} as an independent threshold
construction.

\begin{itemize}[leftmargin=2em,nosep]
\item Pulsar: scheduled for NIST MPTC v0.1 submission on
      $2026$-$11$-$16$ (cut script dry-run validated; see
      \texttt{luxfi/pulsar/scripts/cut-submission.sh}).
\item Corona: scheduled for NIST MPTC v0.2 (target: $2027$ Q1).
\item Magnetar: scheduled for NIST MPTC v0.3 (target: $2027$ Q3).
\end{itemize}

NIST MPTC submission requires Class N1 (single-party-compatible
threshold signing) and Class N4 (public-key preservation across
resharing). Pulsar satisfies both N1 and N4 by construction
(\texttt{TestN1\_ByteEquality} and \texttt{TestN4\_Resharing} in
\texttt{luxfi/pulsar/threshold}); Corona and Magnetar will follow
the same template.

\subsection{Audit status}

The QuasarCert composition (this paper) has been internally
reviewed by the Lux cryptography team. Each individual primitive
has its own audit trajectory:

\begin{itemize}[leftmargin=2em,nosep]
\item Pulsar: internal cryptographer sign-off APPROVED WITH GATES
      ($2026$-$05$-$18$). External audit gated on GATE-4 closure
      (target: 2026 Q4). EC theory shells $13/13$, Lean $5/5$,
      Jasmin $3/3$ at \texttt{v1.0.7}+.
\item Corona: internal cryptographer verdict APPROVED WITH GATES.
      External audit gated on GATE-4 closure (target: 2027 Q1).
\item Magnetar: Tier B; not yet externally audited. The v0.1
      construction has been internally red-team reviewed; the
      \texttt{BLOCKERS.md} file lists open issues.
\item BLS12-381: existing literature audit
      \cite{boneh2001short,boldyreva2003,bonehshorth}; IETF draft
      stabilization in progress.
\item Z-Chain Groth16: Bowe-Gabizon-Miers ceremony audit
      \cite{bowegabizonmiers}; circuit audit included in the
      \texttt{luxfi/chains/quantumvm} repo audit trajectory.
\end{itemize}

The umbrella composition proof (Theorem~\ref{thm:composition}) is
machine-checked in Lean at \texttt{luxfi/proofs/lean/Consensus/Quasar.lean}
\cite{quasarcertsoundness}. The Lean proof tree mirrors the LaTeX
theorem numbers $1$-to-$1$.
